\section{附录：转写校正与复习路线}

本期无平台字幕，自动转写对模型名和英文术语有若干误听。整理时需要把转写中的近似音统一到标准术语，否则会影响后续检索和跨期比较。

\subsection{术语消化：常见误听校正}

\noindent\begin{longtable}{p{0.24\textwidth}p{0.31\textwidth}p{0.36\textwidth}}
\toprule
标准术语 & 常见误听/近似写法 & 校正理由 \\
\midrule
Chatbot & Charbot、拆的 GTT & 第一幕的交互形态。 \\
Anthropic & AnswerPick、Ansropic & 硅谷御三家之一，本期重点公司。 \\
Claude Code & Cloud Code、Cowl Code & Coding/Agent 工作流代表工具。 \\
Codex & CodeX、代码模型泛称 & OpenAI coding 工具体系。 \\
Gemini & Germline、Gemline & Google 前沿模型体系。 \\
Harness Engineering & Harness、哈尼斯工程 & 模型进入真实工作流的中间层。 \\
Model as OS & 模型操作系统、模型 OS & 本期关于未来模型形态的核心判断。 \\
White-collar deflation & 白领通缩 & 社会影响部分的关键词。 \\
\bottomrule
\end{longtable}

\subsection{复习路线}

如果只读一遍，建议按四步复习。第一，读第 2--3 章，理解 Coding 为什么从应用变成加速器。第二，读第 4--6 章，理解模型公司路径依赖、harness 和 model OS。第三，读第 7--10 章，理解社会冲击、投资框架和岗位重组。第四，把本期与 EP139、EP138、EP137 连起来，看 Agent 技术史、模型实验室、AI for Math 和产业季报之间的关系。

\begin{knowledgebox}{本期的最短总结}
Coding 是数字世界的执行语言；Agent 是模型进入环境的方式；Harness 是模型进入工作流的桥；Model OS 是模型成为任务入口的终局想象；白领通缩是生产力跃迁后的社会压力。
\end{knowledgebox}

\subsection{本章小结}

转写校正让术语稳定，复习路线让季度判断可复用。后续继续整理Zhang Xiaojun AI 队列时，本期可以作为“产业季报/综述类”模板。

